\documentclass[11pt]{article}
\usepackage[margin=1in]{geometry}
\usepackage[T1]{fontenc}
\usepackage{lmodern}
\usepackage{microtype}
\usepackage{booktabs}
\usepackage{tabularx}
\usepackage{array}
\usepackage{enumitem}
\usepackage[hidelinks]{hyperref}

\title{LLM-Generated Synthesis Constraints\\Current Project Status}
\author{Alessio Timofte}
\date{June 2026}

\begin{document}
\maketitle

\section*{Research Objective}
The core question is:
\begin{quote}
Can a (local) LLM generate synthesis constraints that are correct enough to support ASIC-style synthesis and timing analysis, and how often do those constraints preserve useful QoR relative to a reference baseline?
\end{quote}

The framework currently studies this question through four linked goals:
\begin{enumerate}[leftmargin=*,nosep]
    \item generate SDC constraints from RTL plus a target clock period,
    \item validate the constraints through a real synthesis and STA flow,
    \item classify the resulting behavior into meaningful failure or success labels,
    \item compare LLM-generated runs against trusted reference constraints.
\end{enumerate}

\section*{Current Architecture}
The repository is now organized around a dedicated Python package, \texttt{pipeline/}, which implements the full experiment loop:

\begin{enumerate}[leftmargin=*,nosep]
    \item \textbf{LLM generation} via \texttt{pipeline/llm.py}. A versioned prompt template (\texttt{v1\_base} or \texttt{v2\_base}, selected at runtime) is filled with Verilog source and target period, then sent to a local Ollama model.
    \item \textbf{Output cleaning} removes reasoning tags, code fences, and non-SDC text, keeping only recognized constraint commands.
    \item \textbf{Synthesis} via \texttt{pipeline/synthesis.py}. Yosys reads the RTL, runs standard optimization and mapping passes, and invokes ABC with a liberty-based timing target.
    \item \textbf{Static Timing Analysis} via \texttt{pipeline/sta.py}. OpenSTA reads the synthesized netlist and the generated SDC, then reports WNS, TNS, minimum slack, and path-level timing behavior. The STA stage issues \emph{independent per-path-group queries} (input$\to$register, register$\to$output, register$\to$register) using \texttt{-from}/\texttt{-to} filters, so a constrained path can no longer be masked by a competing critical path at the same endpoint. This makes coverage detection reliable.
    \item \textbf{Classification} via \texttt{pipeline/classifier.py}. Each run is labeled using both tool logs and the content of the generated SDC.
    \item \textbf{Dataset logging} via \texttt{pipeline/dataset.py}. Each experiment appends one row to \texttt{results/dataset.csv}.
\end{enumerate}

A correction loop is implemented (synthesis-failure trigger); extending it to fire on \texttt{partial\_coverage} and \texttt{timing\_violation} is the primary planned next step.

The framework also supports \emph{multiple prompt versions} and a \texttt{-{}-compare-prompts} mode that runs every template on each design/seed so their results can be compared directly. The complete dataset can be cleanly rebuilt in one invocation using \texttt{-{}-fresh} (drops the existing CSV before the run) combined with \texttt{-{}-with-reference} (regenerates the reference baseline in the same pass).

A separate utility script, \texttt{scripts/rerun\_sta.py}, replays OpenSTA on every existing run's saved netlist and SDC using the current TCL — without re-calling the LLM or Yosys. This is used whenever the STA stage or coverage parser is changed, to upgrade the dataset deterministically ($\sim$6\,s for 60 runs) rather than re-sampling the LLM.

\section*{Toolchain and Environment}
\begin{table}[h]
\centering
\begin{tabularx}{\textwidth}{>{\raggedright\arraybackslash}p{0.24\textwidth} >{\raggedright\arraybackslash}p{0.26\textwidth} X}
\toprule
\textbf{Component} & \textbf{Current Choice} & \textbf{Role in the Project} \\
\midrule
Language model & Ollama + \texttt{qwen3:8b} & Local generation of SDC constraints and correction attempts \\
Synthesis engine & Yosys 0.63 & RTL elaboration, optimization, mapping, reporting \\
Technology mapper & ABC (via Yosys) & Liberty-aware mapping with a delay target derived from \texttt{create\_clock} \\
Timing analyzer & OpenSTA 3.1.0 & Post-synthesis timing analysis and slack extraction \\
Cell library & Nangate45 liberty & Common timing library shared by synthesis and STA \\
Implementation language & Python 3.11+ & Orchestration, parsing, classification, dataset logging \\
\bottomrule
\end{tabularx}
\end{table}

The project remains intentionally lightweight: the only declared Python dependency is the \texttt{ollama} client package, while Yosys and OpenSTA are invoked as subprocess tools. This keeps the framework portable and easy to reproduce on a local macOS setup.

\section*{Benchmark Set}
The current benchmark suite contains six RTL designs, each paired with a trusted reference SDC under \texttt{designs/reference/}:

\begin{table}[h]
\centering
\begin{tabularx}{\textwidth}{>{\ttfamily\raggedright\arraybackslash}p{0.23\textwidth} X}
\toprule
\textbf{Design} & \textbf{Purpose} \\
\midrule
simple.v & Minimal sanity-check design for verifying the end-to-end flow \\
counter.v & Sequential logic with reset and enable behavior \\
fsm.v & Small finite-state machine benchmark \\
adder.v & Registered arithmetic datapath \\
alu.v & Multi-operation combinational and sequential datapath case \\
shift\_reg.v & Serial-in / parallel-out sequential structure \\
\bottomrule
\end{tabularx}
\end{table}

This is a sensible first benchmark layer: the designs are diverse enough to expose clocking, port coverage, and timing-path issues, while still being small enough for fast iterative experimentation.

\section*{What the Framework Measures}
Each run records both correctness and QoR-related information. The current CSV schema captures:
\begin{itemize}[leftmargin=*,nosep]
    \item metadata: design name, top module, model name, seed, prompt version, run directory;
    \item structural properties of the produced SDC: number of extracted commands, presence of \texttt{create\_clock}, number of input and output delay statements;
    \item synthesis metrics: success status, runtime, total cell count, chip area;
    \item timing metrics: STA success, WNS, TNS, minimum slack, timing-met flag, number of setup violations, no-paths detection;
    \item \textbf{per-path-group coverage}: path counts and worst-slack for the in2reg, reg2out, and reg2reg groups, plus a \texttt{coverage\_complete} flag indicating whether the required input/output timing was actually analyzed;
    \item semantic classification: primary label, all labels, and a short rationale.
\end{itemize}

This is a strong design choice because it allows the project to separate several qualitatively different outcomes:
\begin{itemize}[leftmargin=*,nosep]
    \item fully accepted constraints (full path-group coverage, timing met),
    \item \emph{partial-coverage} constraints that pass timing only on the subset of paths that were analyzed (e.g. register-to-register only),
    \item syntactically accepted but ineffective constraints (no timing paths at all),
    \item timing-violating constraints,
    \item synthesis or STA failures caused by bad constraint content.
\end{itemize}

\section*{Current Dataset}
As of \textbf{June 2026}, \texttt{results/dataset.csv} contains 60 runs: 6 reference, 36 LLM runs with the original prompt (\texttt{v1\_base}), and 18 with the revised prompt (\texttt{v2\_base}), all using \texttt{qwen3:8b}. The six reference runs complete successfully with full path-group coverage and serve as the control baseline.

The central result is the comparison between the two prompt versions:

\begin{table}[h]
\centering
\begin{tabular}{lrr}
\toprule
\textbf{Final status} & \textbf{v1\_base (36)} & \textbf{v2\_base (18)} \\
\midrule
ok (full coverage, timing met) & 0 & 18 \\
partial\_coverage              & 20 & 0 \\
ineffective\_no\_paths         & 9 & 0 \\
timing\_violation              & 3 & 0 \\
sta\_failed                    & 4 & 0 \\
\midrule
\textbf{success rate} & \textbf{0\%} & \textbf{100\%} \\
\bottomrule
\end{tabular}
\end{table}

Under \texttt{v1\_base}, \emph{not one} run achieved full coverage. The dominant failure was \emph{partial\_coverage}: the tools ran and reported clean timing, but only on register-to-register paths, because the generated SDC omitted the \texttt{-clock} association required for \texttt{set\_input\_delay} / \texttt{set\_output\_delay} to take effect. Since register-to-register paths are timed by \texttt{create\_clock} alone, they essentially cannot fail --- so naive ``tool acceptance'' badly overstated success. This is the project's key methodological finding: \textbf{tool acceptance alone is not a valid success criterion; path-group coverage must be verified}.

The \texttt{v2\_base} prompt adds exactly this one requirement: name the clock with \texttt{-name}, and attach \texttt{-clock} to every I/O delay. That single change lifted the success rate from 0\% to 100\%. The per-design breakdown confirms the failures were prompt-driven, not design-driven: every design that failed under \texttt{v1\_base} (including \texttt{simple} and \texttt{adder}) is fully accepted under \texttt{v2\_base}.

\paragraph{QoR is fixed by construction.} A netlist-invariance check confirms cell count and chip area are identical to the reference for all runs. This is expected: synthesis is driven only by the clock period, so the SDC's I/O delays cannot change the circuit --- only its timing analysis. Per-group worst-slack comparisons show \texttt{v2\_base} reproduces the reference's slacks exactly (e.g. \texttt{adder} in2reg $+3.72$, reg2out $+4.42$ for both).

\section*{Next Steps}
With the coverage methodology now solid and the prompt-comparison result established, the most valuable next steps are:
\begin{enumerate}[leftmargin=*,nosep]
    \item extend the correction loop to fire on \texttt{partial\_coverage} and \texttt{timing\_violation}, not just synthesis failure, and quantify whether retry-based self-repair lifts \texttt{v1\_base} toward \texttt{v2\_base};
    \item expand the benchmark set with designs that stress multiple clocks, generated clocks, and more complex timing interfaces;
    \item compare multiple Ollama-served models under the same prompt;
    \item explore SDC-aware synthesis (feeding constraints into the mapper) so the more demanding question --- can the LLM \emph{safely relax} timing to reduce area without breaking function --- becomes measurable;
    \item add plots for success rate, coverage, and per-group slack distributions.
\end{enumerate}

\end{document}
